\documentclass[10pt,a4paper]{amsart}
\usepackage[margin=19mm]{geometry}
\usepackage{amsmath,amssymb,mathtools}
\usepackage[scr=rsfs,cal=euler]{mathalfa}
\usepackage{xfrac}
\usepackage[unicode,hidelinks,bookmarks=false]{hyperref}
\newcommand{\ie}{i.e.\ }
\newcommand{\cf}{cf.\ }
\newcommand{\resp}{respectively\ }
\newcommand{\Subsection}{\subsection}
\providecommand{\nobreakdash}{\nobreak\hbox{-}}
\setlength{\emergencystretch}{2em}
\multlinegap0pt
\usepackage{mathtools}
\usepackage{mathrsfs}
\usepackage{tikz-cd}
\usetikzlibrary{decorations.markings, arrows.meta}
\usetikzlibrary{fit}
\hypersetup{
  colorlinks=true,
  linkcolor=blue,
  citecolor=blue,
  urlcolor=blue
}
\newcommand{\josef}[1]{\textcolor{purple}{#1}}
\newcommand{\walter}[1]{\textcolor{orange}{#1}}
\newcommand{\val}[1]{\textcolor{blue}{#1}}
\newtheorem{theorem}{Theorem}[section]
\newtheorem{proposition}[theorem]{Proposition}
\newtheorem{lemma}[theorem]{Lemma}
\newtheorem{corollary}[theorem]{Corollary}
\newtheorem{definition}[theorem]{Definition}
\newtheorem{remark}[theorem]{Remark}
\newtheorem*{definition*}{Definition}
\newtheorem*{lemma*}{Lemma}
\newcommand{\R}{\mathbb{R}}
\newcommand{\N}{\mathbb{N}}
\newcommand{\Sig}{\operatorname{Sig}}
\newcommand{\BV}{\operatorname{BV}}
\newcommand{\id}{\mathbf{1}}
\newcommand{\Tens}{T((\R^d))}
\newcommand{\Trunc}[1]{T^{(#1)}(\R^d)}
\newcommand{\proj}[1]{\pi_{#1}}
\newcommand{\wind}{\operatorname{Wind}}
\newcommand{\Area}{\operatorname{Area}}
\begin{document}
\title[An Elementary Proof of the Hambly-Lyons Uniqueness Theorem]{An Elementary Proof of the Hambly-Lyons Uniqueness Theorem}
\author{Walter Schachermayer, Josef Teichmann, Valentin Tissot-Daguette}
\date{}
\maketitle
\noindent\emph{Source and licence.} An Elementary Proof of the Hambly-Lyons Uniqueness Theorem. Version arXiv:2609.12283v2 (CC BY 4.0). This material is distributed under the Creative Commons Attribution 4.0 International licence, http://creativecommons.org/licenses/by/4.0/, as recorded in the arXiv OAI licence metadata for this exact version and shown on the abstract page https://arxiv.org/abs/2609.12283v2. Full rights notice: licenses/arxiv-2609.12283-CC-BY-4.0.txt.\par\smallskip
\noindent\textit{Editorial scope.} Counted unit: the original \texttt{theorem} environment \texttt{thm:main} at source lines 494--502 together with its complete proof at lines 504--521; the delivered document keeps source lines 167--522 so that every referenced label and every citation resolves inside it. Native editable source is the paper's own concatenated TeX; the AMS wrapper is editorial formatting only. No publisher class, upstream script or unrelated artwork is redistributed.
\begin{equation}\label{eq:LevyArea}
      A(\gamma)
  :=
  \frac12\left(\int_0^1 \gamma^1_u\,d\gamma^2_u
  -\int_0^1 \gamma^2_u\,d\gamma^1_u\right) \, ,
\end{equation}
otherwise called \emph{Lévy area}.

We call $ \Sig_{0,1}(\gamma)$ the (total) signature of the curve $\gamma$. In particular, it holds that $\Gamma$ connects $\id$ and the total signature in $T((\R^d))$, i.e., $\Gamma_1=\Sig_{0,1}(\gamma)$. The natural question arises concerning which properties of $\gamma$ are characterized by its total signature. A complete answer has been given in \cite{HL2010}.

We shall use several basic properties of signatures of continuous bounded variation curves starting at $0$:

\begin{proposition}[Chen's identity \cite{C1954}]\label{prop:chen}
For $ 0\leq r, s, t \leq 1$,
\[
  \Sig_{r,t}(\gamma)=\Sig_{r,s}(\gamma)\Sig_{s,t}(\gamma).
\]
In particular, every $\Sig_{s,t}(\gamma)$ is invertible in $T((\mathbb{R}^d))$. For $ s,t \in [0,1]$, we have
\[
  \Sig_{s,t}(\gamma):=\Sig_{t,s}(\gamma)^{-1} \, .
\]
Chen's identity also holds for truncated signatures of all orders.
\end{proposition}

\begin{remark}
Chen's identity follows from the fact that $\Sig_{s,t}(\gamma)$ for $ s, t \in [0,1]$ is the unique evolution associated to the solution of the ordinary differential equation
$$
d x_u = \sum_{i=1}^d x_u e_i d\gamma^i_u, \; \;   x_s = \id  \;,
$$
on $T((\mathbb{R}^d))$, and  therefore satisfies a semigroup property, which is precisely Chen's identity. The same argument applies to truncated signatures.
\end{remark}

Furthermore, the signature is obviously reparametrization-invariant:

\begin{proposition}[reparametrization invariance]\label{prop:reparam}
Fix $ s < t $ in $[0,1]$. Let $\phi:[0,1]\to[s,t]$ be continuous, increasing and onto. Then
\[
  \Sig_{0,1}(\gamma\circ\phi)=\Sig_{s,t}(\gamma) \, .
\]
\end{proposition}

\begin{proposition}\label{prop:SW}
Consider the continuous functions $C(\Gamma([0,1]))$ on the range 
\begin{equation}\label{eq:SigRange}
    \Gamma([0,1]) = \{\Gamma_t, \ t \in [0,1]\}, 
\end{equation}
which is a compact subset of the locally convex space $T((\R^d))$. Then the restrictions of linear functionals, i.e., the restrictions of elements of $T(\mathbb{R}^d)$ on $\Gamma([0,1])$,  form a point separating sub-algebra and are therefore dense in $C(\Gamma([0,1]))$ with respect to the uniform topology. 
\end{proposition}


\begin{proof}
Since products of signature components are linear combinations of signature components by the Leibniz rule, the restriction of linear functionals on $\Gamma([0,1])$ forms an algebra. This sub-algebra of the algebra of continuous functions is furthermore point separating since the dual space $T(\R^d)$ is point separating as well. Therefore, we can conclude by the Stone--Weierstrass theorem. 
\end{proof}

\begin{lemma}[Linear one-form lemma]\label{lem:polynomial-one-form}
Let $ \ell_1, \ell_2 \in T(\mathbb{R}^d) $ be two linear functionals on $T((\mathbb{R}^d))$. Then
\[
  \int_0^1 \ell_1(\Gamma_t)\,d \ell_2(\Gamma_t)
\]
is a finite linear combination of total signature components of $ \gamma $, i.e., there exists $ \ell \in T(\mathbb{R}^d)$ such that the above integral equals $\ell(\Gamma_1)$. Clearly the signature component of the empty word does not appear in this linear combination.
\end{lemma}

\begin{proof}
Each $w$-coordinate of $\Gamma_t$, i.e., $\Gamma_t^w = S^w_{0,t}(\gamma) \in \R$, $w = i_1\cdots i_n$,
%\footnote{\val{[or as Walter suggested:  "i.e., $\pi_N \Gamma_t \in T^{(N)}(\R^d)$," though these would be the layers in the signature tree, not a single integral]}}  
is an iterated integral of $ \gamma $ over $[0,t] $. Since products of iterated integrals are linear combinations of iterated integrals by the Leibniz rule, the integration of an iterated integral against an iterated integral gives again a finite linear combination of iterated integrals of $ \gamma $ over $ [0,1] $. Hence the value of the integral is a finite linear combination of total signature components.
\end{proof}



%\begin{proposition}\label{prop:levyarea_winding}
%Let $d=2$ and $\gamma(0)=\gamma(1)$. Then the Lévy area can be expressed by winding numbers
%\[
%2 \pi \sqrt{-1} A(\gamma) = \int_{\mathbb{C} \setminus \gamma([0,1])} \int_{\gamma} \frac{1}{z-z_0} d z \, \lambda(dz_0). \val{\quad [\lambda(dz_0) \to dz_0, \text{ or define $\lambda$}]}
%4 \pi \sqrt{-1} A(\gamma) = \frac{\partial^2}{\partial\epsilon^2}|_{\epsilon=0} \int_{\mathbb{C} \setminus \epsilon \gamma([0,1])} %\int_{\epsilon\gamma} \frac{1}{z-z_0} d z \, \lambda(dz_0) \, ,
%\]
%\end{proposition}
%\begin{proof}
%The formula holds true since both sides are apparently reparametrization invariant, strongly continuous with respect to the bounded variation norm, real analytic, and $2$-homogeneous.
%\val{[maybe we can invoke the Jordan curve theorem by decomposing the path using self intersections somehow? outside the Jordan curves the Winding number is 0... ]}
%\end{proof}

In the planar case ($d = 2$), the following proposition relates the Lévy area to the winding numbers of $\gamma$, defined as  
$$\textnormal{Wind}_{z_0}(\gamma) = \frac{1}{2 \pi \mathrm{i}} \int_{\gamma} \frac{1}{z-z_0} d z , \quad z_0 \in \mathbb{C} \setminus \gamma([0,1]) \, .$$
\begin{proposition}\label{prop:levyarea_winding}
Let $d=2$ and $\gamma(0)=\gamma(1)$. Then the Lévy area \eqref{eq:LevyArea} can be expressed as  
\begin{equation}\label{eq:LevyvsWinding}
    A(\gamma) = \int_{\mathbb{C} \setminus \gamma([0,1])}  \textnormal{Wind}_{z_0}(\gamma) \, \lambda(dz_0),
\end{equation}
where $\lambda$ denotes the planar Lebesgue measure. 
\end{proposition}
\begin{proof}
This follows from the more general Cufí--Verdera  formula  \cite[Theorem 1]{CV2015}, or can simply be seen directly by Fubini's theorem and the well-known potential-theoretic formula
$$
\frac{1}{\pi} \int_{B_R} \frac{z-z_0}{{|z-z_0|}^2} \lambda(dz_0)  = z
$$
for a disc $B_R \subset \mathbb{C}$ around zero with radius $ R > |z|$.
\end{proof}

\section{Trees and tree-like paths}\label{sec:tree-like}

Recall that a closed curve on a finite tree splits into elementary excursions: whenever the walk enters a branch, it must later leave that branch along the same arc. General compact trees are treated with the same intuition. 

\begin{definition}[Compact tree]
A compact tree is a compact connected metric space $T$ such that any two points $a,b\in T$ are connected by a unique arc, i.e., a unique, injective continuous curve. This arc is denoted by $[a,b]$.
\end{definition}


\begin{definition}[Tree-like path]
A continuous bounded variation path $ \gamma:[0,1]\to\R^d $ is called tree-like if there exist a compact tree $ T $, a continuous loop
\[
  \tau:[0,1]\to T,
  \qquad \tau(0)=\tau(1),
\]
and a continuous map $ \psi:T\to\R^d $ such that
\[
  \gamma=\psi\circ\tau.
\]
In particular, $ \gamma(1)=\gamma(0) $. Analogously, we speak of tree-like paths on $[s,t]$.
\end{definition}

Tree-like paths can also be described by so-called height functions \cite{HL2010}. Here, we adopt an equivalent definition following \cite{BGLY2016}, which better fits our geometric viewpoint. See also \cite[Theorem 5.15]{L2017} for other equivalent definitions of tree-like paths.


\begin{remark}\label{openmap}
If a continuous bounded variation path $ \gamma:[0,1]\to\R^d $ factors over a continuous loop $ \tau:[0,1] \to T$ through $ \gamma= \psi \circ \tau$, then $\psi|_{\tau([0,1])}$ is automatically continuous.
\end{remark}

We now record the elementary two-dimensional lemma used in the subsequent induction over signature levels.

\begin{lemma}[Tree-like planar curves have vanishing Lévy area]\label{lem:2DTreeLikeLemma}
Let $\beta:[0,1]\to\R^2$ be a continuous bounded variation path. If $\beta$ is tree-like, then
\[
  \Sig^{\leq 2}_{0,1}(\beta)=\id \, .
\]
Equivalently, $\beta$ is a loop  with vanishing Lévy area: $A(\beta) = 0$. 
\end{lemma}

\begin{proof}
  Since $\beta$ is tree-like, it factors as $\beta=\psi\circ\tau$, where $\tau:[0,1] \to T$ is a loop in a compact tree. Hence $\beta(1)=\beta(0)$, and the first level of the signature vanishes.

It remains to consider the second level. We now interpret $\beta$ as a curve in $\mathbb{C}$. Let $ z_0 \notin \beta([0,1]) $. By Cauchy's theorem of complex analysis, $\int_\beta \frac{1}{z-z_0} d z=0$, since $\beta$ can be retracted to a point within its own range by being tree-like \cite[Theorem 5.15]{L2017}. On the other hand, we have for all closed, continuous bounded variation curves
\[
2 \pi \textrm{i} A(\beta) = \int_{\mathbb{C} \setminus \beta([0,1])} \int_{\beta} \frac{1}{z-z_0} d z \, \lambda(dz_0) \, ,
\]  
by Proposition \ref{prop:levyarea_winding}. Hence the L\'evy area $A(\beta)$ vanishes.

The symmetric part of the second level is determined by the identity:
\[
   \Sig^{ij}_{0,1}(\beta)+  \Sig^{ji}_{0,1}(\beta)
  =
  \Sig^i_{0,1}(\beta)  \Sig^j_{0,1}(\beta)=0 \, .
\]
Together with the vanishing of the anti-symmetric part, this gives $\Sig^{ij}_{0,1}(\beta)=0$ for all $i,j\in\{1,2\}$. Hence $\Sig^{\leq 2}_{0,1}(\beta)=\id$.
\end{proof}

\begin{proposition}[Tree-likeness of signature curves]\label{prop:lifts-tree-like}
If $ \gamma:[0,1]\to\R^d $ is tree-like, i.e., it factors over a continuous loop $ \tau:[0,1] \to T $ for some compact tree $T$, then for every $N\geq 1$, the truncated signature curve
\[
  t\longmapsto \Sig^{\leq N}_{0,t}(\gamma)
\]
factors over the same loop $\tau: [0,1] \to T$.
\end{proposition}

\begin{proof}
We argue by induction on $N$. For $N=1$ the lifted path is $t\mapsto \gamma(t)-\gamma(0)$. Thus it is tree-like by assumption and factors over the loop $\tau:[0,1] \to T $.

Assume that
\[
  X_t:=\Sig^{\leq N}_{0,t}(\gamma)
\]
is tree-like factoring over the same loop $ \tau:[0,1] \to T$ as $\gamma$ does. Fix $s<t$ such that $\tau(s)=\tau(t)$; then $\tau|_{[s,t]}$ is a loop in $T$, too. The increment path
\[
  Y_u:=\Sig^{\leq N}_{s,u}(\gamma) = \proj{N}(\Gamma_s^{-1} \Gamma_u),
  \qquad u\in[s,t],
\]
is  tree-like  in \(T^{(N)}(\R^d)\), factoring over this loop $ \tau|_{[s,t]}$. We must prove that every component of level $N+1$ from $\Sig_{s,t}(\gamma)$ vanishes, too.

Let $ w=i_1\cdots i_N$ be a word of length $N$, and let $ j\in\{1,\ldots,d\}$ be a letter. Consider the 
two-dimensional coordinate projection of the lifted and restricted loop $Y$ via 
\[
  P_{w,j}:T^{(N)}(\R^d)\to\R^2,
  \qquad
  P_{w,j}(a):=(a^w,a^j).
\]
Since $Y$ is tree-like, the planar bounded variation path
\[
  Z_u:=P_{w,j}(Y_u)
  =\bigl(S^w_{s,u}(\gamma),S^j_{s,u}(\gamma)\bigr)
  =\bigl(S^w_{s,u}(\gamma),\gamma^j(u)-\gamma^j(s)\bigr) \in \R^2
\]
is tree-like, too,  factoring over the continuous loop $ \tau|_{[s,t]}$. By   Lemma~\ref{lem:2DTreeLikeLemma}, 
its second-order signature vanishes. In particular, its $(1,2)$-component is zero:
\[
  0
  =S^{12}_{s,t}(Z)
  =\int_{s<u<v<t} dZ^1_u\,dZ^2_v
  =\int_s^t Z^1_v\,dZ^2_v = S^{wj}_{s,t}(\gamma) \, .
\]
Since $wj$ is an arbitrary word of length $N+1$, all level $ N+1 $ coordinates of $ \Sig_{s,t}(\gamma) $ vanish. Hence the level $ N+1 $ lift has the property that $ \Sig^{\leq N+1}_{s,t}(\gamma) = \id $, i.e., is trivial. This means that $ \Sig^{\leq N+1}_{0,s}(\gamma) = \Sig^{\leq N+1}_{0,t}(\gamma) $ by Chen's identity.

Thus we have proven that if $\tau(s) = \tau(t)$ for $ s < t $, then  $ \Sig^{\leq N+1}_{0,s}(\gamma) = \Sig^{\leq N+1}_{0,t}(\gamma) $. Therefore, $ t\longmapsto \Sig^{\leq N+1}_{0,t}(\gamma) $ factors over the loop $\tau$, too. The associated map $\psi $ is continuous by Remark \ref{openmap}.
\end{proof}

\section{The sub-interval lemma and the construction of the factoring tree}\label{sec:subinterval_lemma}

We next prove the converse direction, i.e., a path with trivial signature is tree-like (see \cite{BGLY2016}, but there in greater generality). The subsequent lemma contains the key idea and states that a path with trivial total signature must contain a nontrivial subpath, i.e., a restriction on a nontrivial sub-interval $[s,t] \subset [0,1]$, where the total signature is trivial, too.

\begin{lemma}[Sub-interval lemma]\label{lem:subinterval}
Let $\gamma:[0,1]\to\R^d$ be continuous of bounded variation, and suppose
\[
  \Sig_{0,1}(\gamma)=\id \, .
\]
Then there exist \(0\leq s<t\leq 1\) with $ t-s < 1$, i.e., a proper sub-interval, such that
\[
  \Sig_{s,t}(\gamma)=\id.
\]
\end{lemma}

\begin{proof}
Assume, by contradiction, that there is no proper sub-interval with trivial signature. Set
\[
  \Gamma_t:=\Sig_{0,t}(\gamma).
\]
Then $\Gamma_0=\Gamma_1 = \id $ in view of  $ \Sig_{0,1}(\gamma)=\id $. Moreover, if $\Gamma_s=\Gamma_t$ for $ s < t$ and $t-s <1$, then by  Chen's identity,
\[
  \Sig_{s,t}(\gamma)=\Gamma_s^{-1}\Gamma_t= \id \, .
\]
By assumption, this cannot happen. Hence $\Gamma$ defines a simple closed loop, i.e., an injective map,
\[
  \Gamma:S^1\to T((\R^d)), 
\]
where $S^1 $ denotes the $1$-sphere in $\R^2$. Since $S^1$ is compact, $\Gamma$ is a homeomorphism onto its image. Let
\[
  h:\Gamma(S^1)\to S^1\subset\R^2
\]
be its inverse.

The coordinate functions of $h=(h_1,h_2)$ are continuous on the compact set $\Gamma(S^1)$. By Proposition \ref{prop:SW}, the two coordinate functions of $h$ can be uniformly approximated 
by the restriction of linear functionals on $\Gamma(S^1)$. Hence there is a continuous linear map
\[
  l:T((\R^d))\to\R^2
\]
such that $l\circ \Gamma: S^1\to \R^2$ 
is uniformly close to the identity map on $S^1$. Choosing the approximation sufficiently close, the planar curve $ l \circ \Gamma $ stays away from the origin and has nonzero winding number:
\[
  \wind_0( l \circ\Gamma)\neq 0.
\]
An illustration is given in Figure~\ref{fig:winding_approximation}. 
The winding number is determined by integrating the angular one-form 
\[
  \omega=\frac{x\,dy-y\,dx}{x^2+y^2}.
\]
Since $ l \circ \Gamma $ remains in a compact subset of $\mathbb{R}^2\setminus\{0\}$, the form $\omega$ can be uniformly approximated there by polynomial one-forms. Applying Proposition \ref{prop:SW}, namely that the linear functionals on $\Gamma(S^1)$ form an algebra, we finally obtain linear functionals $\ell_1^k, \ell_2^k \in T(\mathbb{R}^d)$, such that the approximate winding number can be written as
\[
\sum_k  \int_0^1 \ell_1^k(\Gamma_t) \, d \ell_2^k(\Gamma_t) \neq 0 \, .
\]


By Lemma~\ref{lem:polynomial-one-form}, these integrals, however, are a finite linear combination of total signature coordinates of $\gamma$ for nonempty words, hence are determined by $\Sig_{0,1}(\gamma)$. Since $ \Sig_{0,1}(\gamma)=\id $, all signature components for nonempty words vanish, whence the integrals have to vanish, too, which is the desired contradiction.

Therefore, a proper sub-interval $ [s,t] $ with $ \Sig_{s,t}(\gamma)=\id $ must exist.
\end{proof}

\subsection*{Factoring Tree Construction}


Recall the signature range  
$$\Gamma([0,1]) = \{\Gamma_t, t\in [0,1]\} \subset T((\R^d))$$ introduced in Proposition~\ref{prop:SW}. 

\begin{figure}[t]
    \centering
    \begin{tikzpicture}[scale=1.05,
        ->-/.style={postaction={decorate,decoration={
            markings,
            mark=at position #1 with {\arrow[scale=1.5]{stealth}}
        }}}
    ]
        % Origin
        \fill (0,0) circle (2pt) node[below right] {$0$};

        % Unit circle
        \draw[->-={0.135}, thick, black, densely dashed] (0,0) circle (2cm);
        \node[black, right] at (1.75,1.05) {$S^1$};

        % l \circ \Gamma
        \draw[
            ->-={0.383},
            thick,
            black,
            domain=0:360,
            samples=300,
            variable=\t,
            line join=round
        ]
        plot (
            {\t + ((\t>36)*(\t<84)
                 + (\t>180)*(\t<216)
                 + (\t>240)*(\t<288))
                 *10*sin(15*\t)}
            :
            {2 + 0.14*sin(10*\t) + 0.13*cos(15*\t)}
        );

        \node[black, above left] at (-1.8,0.99) {$l \circ \Gamma$};
    \end{tikzpicture}
    \caption{Map $l \circ \Gamma$ (solid) approximating the standard parametrization
    of $S^1$ (dashed) in the proof of Lemma~\ref{lem:subinterval} and preserving
    the non-zero winding number around the origin.}
    \label{fig:winding_approximation}
\end{figure}


\begin{theorem}[Existence of a factoring tree]\label{thm:quotient-tree}
Let $ \gamma:[0,1]\to\R^d $ be a continuous path of bounded variation such that $ \Sig_{0,1}(\gamma)=\id $. Then, 
$$ T_\gamma := \Gamma([0,1]) \subset T((\R^d))$$
is a compact tree. Moreover, there is a continuous map $  \psi:T_\gamma\to\R^d$ such that 
  $\gamma=\psi\circ \Gamma$. 
\end{theorem}
\begin{proof}
Let $a,b \in T_\gamma$ with $a \neq b$. We have to show that there is a unique (up to reparametrization) continuous injective arc connecting $a$ to $b$. If this were not the case, there would exist two distinct arcs $\alpha$ and $\beta$ connecting $a$ with $b$. Since they are different, there are proper subarcs $\alpha|_{[u_-,u_+]}$ and $\beta|_{[v_-,v_+]}$ such that $\alpha(u_-)=\beta(v_-)$, $\alpha(u_+)=\beta(v_+)$, but otherwise disjoint. Concatenating those subarcs leads to an injective loop in $T_\gamma$, which contradicts the sub-interval lemma. Therefore, $T_\gamma$ is a compact tree. Hence $\gamma$ is tree-like. 

 For the second assertion, we may assume without loss of generality that $\gamma(0) = 0$. We then define $\psi:T((\R^d))\to\R^d$ as the canonical projection of the signature onto its first level, i.e.,  the curve itself. 
\end{proof}

\section{Hambly-Lyons uniqueness theorem}


\begin{theorem}[Hambly--Lyons uniqueness]\label{thm:main}
Let $ \gamma:[0,1]\to\R^d $ be a continuous bounded variation path. Then
\[
  \Sig_{0,1}(\gamma)=\id
  \quad\Longleftrightarrow\quad
  \gamma \text{ is tree-like}.
\]
Consequently, the signature separates based continuous bounded variation paths up to tree-like equivalence, in the sense of Hambly--Lyons \cite{HL2010}.
\end{theorem}

\begin{proof}
Assume first that $\gamma$ is tree-like. By Proposition~\ref{prop:lifts-tree-like}, for every $N\geq1$,  the lifted path
\[
  t\longmapsto \Sig^{\leq N}_{0,t}(\gamma)
\]
is tree-like and closed. Hence its endpoint agrees with its starting point, and therefore
\[
  \Sig^{\leq N}_{0,1}(\gamma)=\id
\]
for every $N$. Thus $\Sig_{0,1}(\gamma)=\id$.

Conversely, assume $\Sig_{0,1}(\gamma)=\id$. By Theorem~\ref{thm:quotient-tree}, $\gamma$ factors as
\[
  \gamma=\psi \circ \Gamma, 
\]
where $T_\gamma =  \Gamma([0,1])$ is a compact tree. 
Hence $\gamma$ is tree-like.
\end{proof}

\begin{thebibliography}{99}
\bibitem[BGLY2016]{BGLY2016} Horatio Boedihardjo, Xi~Geng, Terry Lyons, and Danyu Yang. \newblock The signature of a rough path: Uniqueness. \newblock {\em Advances in Mathematics}, 293:720--737, 2016.
\bibitem[C1954]{C1954} Kuo-Tsai Chen, \newblock Iterated integrals and exponential homomorphisms, \newblock {\em Proceedings of the London Mathematical Society} \textbf{s3-4} (1954), no.~1, 502--512.
\bibitem[CV2015]{CV2015} Juli{\`a} Cuf{\'\i} and Joan Verdera. \newblock A general form of {G}reen's formula and the {C}auchy integral theorem. \newblock {\em Proceedings of the American Mathematical Society}, 143(5):2091--2102, 2015.
\bibitem[HL2010]{HL2010} B.~M. Hambly and T.~J. Lyons. \newblock Uniqueness for the signature of a path of bounded variation and the reduced path group. \newblock \emph{Annals of Mathematics}, 171(1):109--167, 2010.
\bibitem[L2017]{L2017} Thierry L{\'e}vy. \newblock {\em The master field on the plane}. \newblock {\em Ast{\'e}risque}, 388:ix+201, 2017.
\end{thebibliography}
\end{document}
